\documentclass[11pt]{article}
\usepackage[margin=25mm]{geometry}
\usepackage{amsmath,amssymb,amsthm}
\usepackage{hyperref}
\newtheorem{lemma}{Lemma}
\newtheorem{theorem}{Theorem}
\newtheorem{proposition}{Proposition}
\newcommand{\R}{\mathbb R}
\newcommand{\one}{\mathbf 1}
\newcommand{\sig}{\sigma}
\title{Stationary AR(1) information under missing observations:\\ exact identities and the full-history profiling boundary}
\author{Research working note, version 1}
\date{7 October 2026}
\begin{document}
\maketitle

\paragraph{Scope.}
The Gaussian, Markov and Schur identities below are tools, not claims of a
new control-theoretic contribution. The sharp minimax coefficient conjecture
is not proved or disproved by this note. The companion Chinese note gives
the research contract and remaining obligations.

Let $n\ge1$, $\sig>0$, $|\rho|<1$, and let stationary Gaussian noise satisfy
$\epsilon_i=\rho\epsilon_{i-1}+\xi_i$, where
$\epsilon_1\sim N(0,\sig^2)$ and the independent innovations have variance
$\sig^2(1-\rho^2)$. For a deterministic retained set
$R=\{t_1<\cdots<t_p\}\subset\{1,\ldots,n\}$, its marginal covariance is
$C_R=(\sig^2\rho^{|t_i-t_j|})$. For any deterministic mean contrast $x$ write
\begin{equation}\label{eq:contract}
 Q_R(x)=x_R^\top C_R^{-1}x_R,\qquad G_R(x)=\tfrac12Q_R(x).
\end{equation}
Here $G$ is Gaussian pairwise KL and $Q$ is squared Mahalanobis distance;
both are zero for $R=\varnothing$. Define Fisher (not half-KL) coefficients
\begin{equation}\label{eq:coefficients}
 w_d=\frac{1-\rho^d}{1+\rho^d},\quad
 \nu=\frac{w_1}{\sig^2},\quad
 \beta(k)=\frac{(k+1)w_1-w_{k+1}}{\sig^2}.
\end{equation}

\begin{lemma}[Marginal energy]\label{lem:marginal}
For $p\ge1$, $d_j=t_j-t_{j-1}$,
\begin{equation}\label{eq:marginal}
 Q_R(x)=\frac{x_{t_1}^2}{\sig^2}
 +\sum_{j=2}^p\frac{(x_{t_j}-\rho^{d_j}x_{t_{j-1}})^2}
 {\sig^2(1-\rho^{2d_j})}.
\end{equation}
\end{lemma}
\begin{proof}
Iteration gives
$\epsilon_{i+d}=\rho^d\epsilon_i+\sum_{h=1}^d\rho^{d-h}\xi_{i+h}$.
The second term is independent of the noise history through $i$, and has
variance $\sig^2(1-\rho^2)\sum_{h=0}^{d-1}\rho^{2h}
=\sig^2(1-\rho^{2d})$. Innovation sums over disjoint retained intervals are
independent. Thus the marginal density factors into one stationary initial
density and these transition densities. Their triangular whitening maps the
mean contrast to $x_{t_1}/\sig$ and
$(x_{t_j}-\rho^{d_j}x_{t_{j-1}})/(\sig\sqrt{1-\rho^{2d_j}})$.
The squared norm of this vector is $Q_R$, proving the result, including
negative $\rho$ since all variances remain positive.
\end{proof}

\begin{lemma}[Constant contrasts, all boundary cases]\label{lem:constant}
If $x=a\one$, deleting one internal gap of $k\ge0$ samples with its endpoints
retained loses $a^2\beta(k)$ in $Q$, or $a^2\beta(k)/2$ in KL. More generally,
for nonempty $R$, set $p_0=t_1-1$, $p_1=n-t_p$ and
$k_j=t_{j+1}-t_j-1$. Then
\begin{equation}\label{eq:multi}
 Q_{[n]}(a\one)-Q_R(a\one)
 =a^2\left\{\nu(p_0+p_1)+\sum_{j=1}^{p-1}\beta(k_j)\right\}.
\end{equation}
For empty $R$ the loss is $a^2[\sig^{-2}+(n-1)\nu]$.
Moreover $\beta(0)=0$, $\beta(k)>0$ for $k\ge1$, and at $\rho=0$ it equals
$k/\sig^2$. No separation condition on distinct gaps is needed.
\end{lemma}
\begin{proof}
A $d$-step term in \eqref{eq:marginal} at constant contrast is
$a^2(1-\rho^d)^2/[\sig^2(1-\rho^{2d})]=a^2w_d/\sig^2$.
The full record therefore has energy $a^2[\sig^{-2}+(n-1)\nu]$.
Replacing $m=k+1$ one-step transitions by one $m$-step transition gives
$a^2(mw_1-w_m)/\sig^2$. Since
$n-1=p_0+p_1+\sum_j(k_j+1)$, subtraction proves \eqref{eq:multi}.
The empty case follows from zero observed information.

For strict positivity, let $P$ be the precision of the local complete
$m+1$ observations; $M$ indexes the $k$ missing interior entries and $E$ the
two endpoints. Completing the square gives
\begin{equation}\label{eq:schur}
 Q_{\rm local}(v)-Q_E(v_E)
 =(v_M+P_{MM}^{-1}P_{ME}v_E)^\top P_{MM}
 (v_M+P_{MM}^{-1}P_{ME}v_E).
\end{equation}
Here $P_{MM}$ is positive definite. Every interior row of $P$ applied to
$a\one$ gives
$a(1-\rho)^2/[\sig^2(1-\rho^2)]=a\nu\ne0$ for $a\ne0$.
Thus the right side is strictly positive whenever $k\ge1$.
The $k=0$ and $\rho=0$ claims follow directly from \eqref{eq:coefficients}.
\end{proof}

\begin{lemma}[Local loss for any contrast]\label{lem:local}
For an internal gap from $\ell$ to $\ell+m$, the exact quadratic loss is
\begin{equation}\label{eq:loss}
 L_{\ell,m}(x)=
 \sum_{h=1}^m\frac{(x_{\ell+h}-\rho x_{\ell+h-1})^2}{\sig^2(1-\rho^2)}
 -\frac{(x_{\ell+m}-\rho^m x_\ell)^2}{\sig^2(1-\rho^{2m})}.
\end{equation}
This is positive semidefinite and zero when $m=1$.
For nonempty $R$, total loss is the sum of these internal losses, the prefix
\begin{equation}\label{eq:prefix}
 \frac{x_1^2-x_{t_1}^2}{\sig^2}
 +\sum_{i=2}^{t_1}\frac{(x_i-\rho x_{i-1})^2}{\sig^2(1-\rho^2)},
\end{equation}
and the suffix sum over $i=t_p+1,\ldots,n$ of the same one-step terms.
\end{lemma}
\begin{proof}
Subtract \eqref{eq:marginal} for retained data from its full-record version.
Every unaffected one-step term cancels; each retained-to-retained gap leaves
\eqref{eq:loss}. The initial and final parts give \eqref{eq:prefix} and the
suffix. Positivity of each local loss follows by completing the square as in
\eqref{eq:schur} for its local missing subset. For $m=1$ the two terms are
identical. This proof permits gaps sharing a single retained endpoint.
\end{proof}

Put $c_\rho=\rho/[\sig^2(1+\rho)]$ and
$d_\rho=\rho/[\sig^2(1-\rho^2)]$.
\begin{lemma}[Retained-block boundary identity]\label{lem:block}
For each internal missing gap with span $m\ge2$, let its retained endpoints
have contrast $u,v$ and set
\begin{equation}\label{eq:boundary}
 B_m(u,v)=\left[\frac1{\sig^2(1-\rho^{2m})}
 -\frac1{\sig^2(1+\rho)}\right](u^2+v^2)
 -\frac{2\rho^muv}{\sig^2(1-\rho^{2m})}.
\end{equation}
For every nonempty retained set,
\begin{equation}\label{eq:block}
 Q_R(x)=\nu\sum_{i\in R}x_i^2+c_\rho(x_{t_1}^2+x_{t_p}^2)
 +d_\rho\sum_{i,i+1\in R}(x_{i+1}-x_i)^2+\sum_{\rm gaps}B_m(u,v).
\end{equation}
The individual boundary terms need not be nonnegative for negative $\rho$.
\end{lemma}
\begin{proof}
Expanding \eqref{eq:marginal} on a contiguous interval $I$ gives
\begin{equation}\label{eq:contiguous}
 Q_I(x)=\nu\sum_{i\in I}x_i^2+c_\rho(x_{\min I}^2+x_{\max I}^2)
 +d_\rho\sum_{i,i+1\in I}(x_{i+1}-x_i)^2.
\end{equation}
Indeed, its interior diagonal coefficient is
$\nu+2d_\rho=(1+\rho^2)/[\sig^2(1-\rho^2)]$, its endpoint coefficient is
$\nu+d_\rho+c_\rho=1/[\sig^2(1-\rho^2)]$, and its adjacent cross coefficient
is $-2d_\rho$. For a singleton $\nu+2c_\rho=\sig^{-2}$, so the identity
also covers singleton blocks. Sum \eqref{eq:contiguous} over maximal retained
blocks. Every later block's artificial initial term $v^2/\sig^2$ must be
replaced by $(v-\rho^m u)^2/[\sig^2(1-\rho^{2m})]$. Together with its two
internal endpoint corrections this yields
$c_\rho u^2+(c_\rho-\sig^{-2})v^2+
(v-\rho^m u)^2/[\sig^2(1-\rho^{2m})]$, which expands to
\eqref{eq:boundary}. The first and last corrections remain explicit.
\end{proof}
In particular,
\begin{equation}\label{eq:sign}
 B_m(A,-A)-B_m(A,A)=\frac{4\rho^m A^2}{\sig^2(1-\rho^{2m})}.
\end{equation}

\begin{lemma}[Near-constant local bound]\label{lem:smooth}
Suppose $|x_{\ell+h}-x_\ell|\le Lh$ for $h=0,\ldots,m$ and let $A=x_\ell$.
Set
$C_\rho=(1+|\rho|)/[\sig^2(1-|\rho|)]$,
$S_m=m(m+1)(2m+1)/6$ and $D_m=C_\rho L^2S_m$. Then
\begin{equation}\label{eq:smooth}
 |L_{\ell,m}(x)-\beta(m-1)A^2|
 \le2|A|\sqrt{\beta(m-1)D_m}+D_m.
\end{equation}
\end{lemma}
\begin{proof}
The local loss matrix $M$ from \eqref{eq:schur} satisfies $0\preceq M\preceq P$.
The full local precision has absolute interior row sum $C_\rho$ and endpoint
row sum $1/[\sig^2(1-|\rho|)]$, hence norm at most $C_\rho$.
For $m=1$ the loss is zero. Otherwise write $x=A\one+\delta$.
Then $\delta^\top M\delta\le C_\rho\sum_h\delta_h^2\le D_m$ and
$\one^\top M\one=\beta(m-1)$ by Lemma~\ref{lem:constant}.
Positive semidefinite Cauchy--Schwarz gives
$|\one^\top M\delta|\le\sqrt{\beta(m-1)D_m}$.
Expanding $x^\top Mx$ and bounding its cross and remainder terms proves
\eqref{eq:smooth}.
\end{proof}

\begin{theorem}[Conditional full-history profiling bridge]\label{thm:bridge}
Fix $K\ge1$, $A_0,L\ge0$. Let a nonempty full-history class $X_n\subset\R^n$
satisfy for every $x\in X_n$,
$|x_i|\le A_0n$ and $|x_{i+1}-x_i|\le L$.
Assume $1,n\in R$, and every internal gap length $k_j$ is between $1$ and $K$.
There are $J$ such gaps, with left endpoints $\ell_j$. Define
\begin{equation}\label{eq:surrogate}
 T_R(x)=\nu\sum_{i=1}^n x_i^2-\sum_{j=1}^J\beta(k_j)x_{\ell_j}^2,
 \quad \beta_*=(K+1)\nu,\quad D_*=C_\rho L^2S_{K+1}.
\end{equation}
Uniformly over all these calendars and all $x\in X_n$,
\begin{align}\label{eq:uniform}
 |Q_R(x)-T_R(x)|&\le E_n,\\
 E_n&=2|c_\rho|A_0^2n^2+|d_\rho|(n-1)L^2
 +J\{2A_0n\sqrt{\beta_*D_*}+D_*\}.\nonumber
\end{align}
Consequently,
\begin{equation}\label{eq:inf}
 \left|\inf_{x\in X_n}\tfrac12Q_R(x)-\inf_{x\in X_n}\tfrac12T_R(x)\right|
 \le\tfrac12E_n=O(n^2)=o(n^{5/2}).
\end{equation}
The bound is uniform over $|\rho|\le q<1$, $\sig\ge\sig_{\min}>0$ for fixed
$K,A_0,L$.
\end{theorem}
\begin{proof}
Use \eqref{eq:contiguous} on the full interval. Its difference from
$\nu\sum_i x_i^2$ is bounded by
$2|c_\rho|A_0^2n^2+|d_\rho|(n-1)L^2$.
Since both ends are retained, Lemma~\ref{lem:local} subtracts only the internal
losses. Apply Lemma~\ref{lem:smooth} to each. Because $w_m>0$,
$0\le\beta(k_j)=m_j\nu-w_{m_j}/\sig^2\le(K+1)\nu=\beta_*$,
and $D_{m_j}\le D_*$. Summation proves \eqref{eq:uniform}.
The pointwise bounds $T_R-E_n\le Q_R\le T_R+E_n$, when infimized over the
same nonempty $X_n$, prove \eqref{eq:inf}; no attainment or exchange of
blockwise infima is required. Finally $J\le n$ and all displayed constants
are bounded on the stated channel class.
\end{proof}

\paragraph{Nuisance quantifiers.}
For the raw model $Y_i=s_i b_i+\mathbf1_{\rm alt}a_i+\epsilon_i$, define the
actual difference image
$\mathcal Z_n=\{b_0-b_1:b_0\in\mathcal B_{0,n},b_1\in\mathcal B_{1,n}\}$.
Profiled pairwise KL is
$\inf_{z\in\mathcal Z_n}Q_R(a-sz)/2$, always using a single full-history pair.
A Lipschitz constraint on $z$ does not imply the hypothesis of
Theorem~\ref{thm:bridge} for $x=a-sz$. Neither does an unbounded-amplitude
nuisance class imply its envelope hypothesis. Restricting the class requires
a separate justified argument. The geometric infimum is not automatically
the minimax testing risk. Under task demodulation $D_s$, covariance must also
become $D_sC_RD_s$.

\begin{theorem}[Raw full-class boundary bridge]\label{thm:raw}
Suppose $|a_{i+1}-a_i|\le\eta$ and every full-history path in the actual
nonempty class $\mathcal Z_n$ satisfies $|z_{i+1}-z_i|\le r$.
No amplitude bound is imposed. Write $x=a-sz$, $s_i\in\{-1,1\}$.
Let $E$ consist of successive retained pairs $(\ell,\ell+m)$ for which
$m\ge2$ or the task signs differ. With $B_m$ defined by
\eqref{eq:boundary}, also for $m=1$, set
\begin{equation}\label{eq:rawsurrogate}
 S_R(z)=\nu\sum_{i\in R}x_i^2+c_\rho(x_{t_1}^2+x_{t_p}^2)
 +\sum_{(\ell,\ell+m)\in E}B_m(x_\ell,x_{\ell+m}).
\end{equation}
For every $z\in\mathcal Z_n$,
\begin{align}\label{eq:rawerror}
 Q_R(x)-S_R(z)
 &=d_\rho\sum_{\substack{i,i+1\in R\\s_i=s_{i+1}}}
 [\Delta a_i-s_i\Delta z_i]^2,\\
 |Q_R(x)-S_R(z)|&\le |d_\rho|(n-1)(\eta+r)^2.\nonumber
\end{align}
Consequently the infima of $Q_R(a-sz)/2$ and $S_R(z)/2$ over this same
full-history class differ by at most $|d_\rho|(n-1)(\eta+r)^2/2=O(n)$.
The bound is uniform for $|\rho|\le q<1$ and $\sig\ge\sig_{\min}>0$.
\end{theorem}
\begin{proof}
The proof of Lemma~\ref{lem:block} permits splitting a contiguous retained
block at an adjacent task change. Substituting $m=1$ into
\eqref{eq:boundary} yields $B_1(u,v)=d_\rho(u-v)^2$.
Thus the only remaining within-block differences in \eqref{eq:block} have
identical task signs. For those edges,
$\Delta x_i=\Delta a_i-s_i\Delta z_i$, proving \eqref{eq:rawerror}.
Every square is at most $(\eta+r)^2$ and there are at most $n-1$ such edges.
Taking infima in the pointwise bounds over the same nonempty class proves
the result. No attainment or amplitude envelope is needed. In particular
$Q_R\ge0$ and the bounded difference ensure $S_R$ has a finite lower bound
on that class, even if its unrestricted matrix is not positive semidefinite.
For $\rho\ge0$, $S_R\le Q_R$; for $\rho<0$ the inequality reverses.
The uniformity follows from the displayed formula for $d_\rho$.
\end{proof}
This surrogate retains all task and missing-observation endpoint products;
it cannot be reduced further to scalar constant-gap penalties without a
new argument. Such endpoint terms can contribute at order $n^{5/2}$.

\begin{lemma}[Affine gap reflection]\label{lem:reflection}
For one internal gap with span $m=k+1$, let $h_j=j-m/2$, $j=0,\ldots,m$.
If its local contrast is $x=A\one+\eta h$, then
\begin{equation}\label{eq:reflection}
 L(x)=\beta(k)A^2+\eta^2 E_k,\qquad E_k\ge0.
\end{equation}
\end{lemma}
\begin{proof}
Let $M$ be the local loss matrix and $J$ the reflection matrix.
Stationary covariance depends only on index distance, hence $JCJ=C$ and
$JC^{-1}J=C^{-1}$. The endpoint marginal precision, embedded into the full
local dimension, also commutes with $J$. Their difference is $M$, so
$JMJ=M$. Since $J\one=\one$ and $Jh=-h$,
$\one^\top Mh=-\one^\top Mh=0$.
Lemma~\ref{lem:constant} gives $\one^\top M\one=\beta(k)$.
Expand the quadratic and set $E_k=h^\top Mh\ge0$ by local loss positivity.
For $k=0$ the loss matrix is zero.
\end{proof}
To apply this identity to a raw zero-endpoint tent construction, $z$ must
vanish on the entire local $m+1$ slots, including both retained endpoints;
vanishing only on missing slots is insufficient.

\begin{proposition}[Failure of an all-path constant-gap bridge]\label{prop:counter}
For $\sig=1$, $\rho=1/2$, $k=1$, there exist legal raw full-history contrasts
with rate-zero nuisance and $J=\sqrt H$ gaps in horizon $n=H$, along square
integers, such that replacing each exact gap loss by
$\beta(1)x_\ell^2$ gives quadratic error $(8/5)H^{5/2}$, or KL error
$(4/5)H^{5/2}$.
\end{proposition}
\begin{proof}
For sufficiently large square $H$, delete the separated internal samples
$g_j=\lfloor jH/(J+1)\rfloor$, $j=1,\ldots,J$. Set $z_i=H$, $a_i=0$ and let
the task sign flip immediately after each missing sample, which retains the
preceding task. Then $x_i=-s_iH$ is a contrast in the raw model, with the same
raw covariance and one full-history rate-zero nuisance. Every gap has local
contrast $(A,A,-A)$ with $A=\pm H$. Its two full transition terms equal
$A^2/3+3A^2$, whereas its retained two-step term equals $5A^2/3$.
Its exact loss is therefore $5H^2/3$. Since
$\beta(1)=2/3-3/5=1/15$, the difference is $8H^2/5$ per gap.
Lemma~\ref{lem:local} adds these differences exactly. Substituting $J=\sqrt H$
and dividing by two proves both asserted errors.
\end{proof}

\paragraph{What this counterexample does not show.}
For zero signal the legal nuisance $z=0$ makes profiled KL zero. Thus the
example disproves an all-path uniform replacement, not the conjectured
coefficient after profiling. If the oracle does not define tasks during
missing samples, the full-record loss needs a separate definition; the
retained-data identity \eqref{eq:block} remains valid.

\paragraph{Computational support.}
The standard-library Fraction program
\texttt{outputs/ar1\_information\_checks\_v1.py} compares covariance inversion,
Markov energy, constant multigap loss and arbitrary-contrast block identity
for all retained subsets of $n=1,\ldots,8$ and
$\rho\in\{-4/5,-1/2,0,1/2,4/5\}$: 2550 exact checks for each comparison,
plus 2550 raw full-class bridge checks, 35 smooth-gap checks,
35 affine-reflection checks and three raw-jump examples. All passed.
The saved report is \texttt{outputs/ar1\_information\_checks\_v1.json}.
Run the script with standard Python from the repository root.
These finite checks validate implementation; the universal claims rely on
the proofs above. No mechanical scan or sharp calendar theorem is claimed.
\paragraph{Compilation status.}
The built-in compiler reported the platform error
\texttt{Unable to find standard directories for platform}, without source
diagnostics. Compilation is unverified; this does not affect the exact Python
checks. No alternative TeX runtime was installed.
\end{document}
